\section{如何使用学到的模型}

\subsection{方法 1：基于模型的规划（Planning）}

有了动力学模型后，可以直接在模型中搜索最优动作序列：

$$\max_{\mathbf{a}_{t:t+H}} \sum_t r(\mathbf{s}_t, \mathbf{a}_t)$$

\begin{figure}[H]
\centering
\includegraphics[width=0.9\textwidth]{slides-images/slide-15.jpg}
\caption{方法 1b：通过采样进行无梯度优化的规划\protect\footnotemark}
\end{figure}
\footnotetext{Slides 第 15 页。}

具体算法流程：
\begin{enumerate}
    \item 运行某个策略（如随机策略）收集数据 $\mathcal{D} = \{(\mathbf{s}, \mathbf{a}, \mathbf{s}')_i\}$
    \item 学习模型 $f_\phi(\mathbf{s}, \mathbf{a})$
    \item 迭代地采样动作序列，通过模型前向推演选择最优动作
\end{enumerate}

\subsubsection{随机打靶法（Random Shooting）}

最简单的方法是随机采样 $N$ 个动作序列，在模型中 rollout 每个序列，选择累积奖励最高的那个序列的第一个动作。

\subsubsection{CEM（Cross-Entropy Method）}

CEM 是一种更高级的无梯度优化方法：
\begin{enumerate}
    \item 初始化动作分布（如高斯分布）
    \item 从分布中采样 $N$ 个动作序列
    \item 在模型中评估每个序列的累积奖励
    \item 选择 top-$K$ 的"精英"序列
    \item 用精英序列更新动作分布的均值和方差
    \item 重复若干轮迭代
\end{enumerate}

\begin{knowledgebox}{MPC（Model Predictive Control）}
在实际执行时，通常采用 MPC 范式：每一步只执行规划得到的第一个动作，然后在新的状态上重新规划。这可以部分缓解模型在长时间跨度上的累积误差。
\end{knowledgebox}

\subsubsection{方法 1a：基于梯度的规划}

如果模型可微，可以直接对动作序列求梯度来优化。但这种方法在实践中容易陷入局部最优，且对模型的光滑性要求较高。

\subsection{方法 2：使用模型生成合成数据}

\begin{figure}[H]
\centering
\includegraphics[width=0.9\textwidth]{slides-images/slide-25.jpg}
\caption{Model-Based 策略优化的两种路线\protect\footnotemark}
\end{figure}
\footnotetext{Slides 第 25 页。}

规划方法的主要局限：
\begin{itemize}
    \item 测试时计算开销大（每步都需要优化）
    \item 受限于短视野问题
\end{itemize}

解决长视野问题的两种思路：
\begin{enumerate}
    \item \textbf{带终端值函数的规划}：用短视野规划 + 学到的值函数作为终端估值
    \item \textbf{Dyna 风格方法}：用模型生成合成数据来增强 model-free RL 的训练
\end{enumerate}

\begin{importantbox}{Dyna 算法}
Dyna 是 model-based RL 的经典范式：
\begin{enumerate}
    \item 收集真实数据，训练/更新动力学模型
    \item 用模型生成合成数据（synthetic rollouts）
    \item 将合成数据与真实数据混合，用 model-free RL 算法更新策略
\end{enumerate}
核心优势是利用模型的泛化能力将有限的真实数据"放大"为大量训练数据。
\end{importantbox}

\subsection{处理模型误差}

\begin{warningbox}{模型误差的累积效应}
模型在每一步的小误差会沿时间轴累积。如果规划视野为 $H$，且每步误差为 $\epsilon$，则 $H$ 步后的总误差可能高达 $O(H \cdot \epsilon)$ 甚至更大（如果动力学是不稳定的）。因此，盲目信任模型会导致策略"利用"模型的缺陷（model exploitation）。
\end{warningbox}

应对策略：
\begin{itemize}
    \item \textbf{模型集成（Ensemble）}：训练多个模型，用不同模型的预测来估计不确定性
    \item \textbf{短视野 rollout}：限制模型 rollout 的步数
    \item \textbf{持续更新模型}：随着收集新数据不断更新模型
\end{itemize}

\subsection{本章小结}

模型可以通过两种方式使用：(1) 直接规划（优化动作序列），(2) 生成合成数据增强 model-free RL。两种方式都需要妥善处理模型误差。

